\documentclass[a4paper,12pt]{article}

% -------------------------------------------------
% Кодировка и язык
% -------------------------------------------------
\usepackage[T2A]{fontenc}
\usepackage[utf8]{inputenc}
\usepackage[russian]{babel}

% -------------------------------------------------
% Математика
% -------------------------------------------------
\usepackage{amsmath,amssymb,amsthm,mathtools}
\usepackage{icomma}

% -------------------------------------------------
% Шрифт и типографика
% -------------------------------------------------
\usepackage{helvet}
\renewcommand{\familydefault}{\sfdefault}
\usepackage{microtype}
\usepackage{setspace}
\onehalfspacing
\usepackage{parskip}

% -------------------------------------------------
% Страница
% -------------------------------------------------
\usepackage[
  a4paper,
  left=19mm,
  right=19mm,
  top=18mm,
  bottom=20mm,
  headheight=15pt
]{geometry}

% -------------------------------------------------
% Графика
% -------------------------------------------------
\usepackage{tikz}
\usetikzlibrary{calc}
\usepackage{tkz-euclide}
\usepackage{pgfplots}
\pgfplotsset{compat=1.18}

% -------------------------------------------------
% Таблицы и списки
% -------------------------------------------------
\usepackage{tabularx,booktabs,longtable}
\usepackage{enumitem}

% -------------------------------------------------
% Заголовки, колонтитулы, ссылки
% -------------------------------------------------
\usepackage{xcolor}
\usepackage{titlesec}
\usepackage{fancyhdr}
\usepackage{hyperref}

% -------------------------------------------------
% Палитра
% -------------------------------------------------
\definecolor{accent}{HTML}{008080}
\definecolor{darkaccent}{HTML}{004D4D}
\definecolor{textgray}{HTML}{333333}
\definecolor{softgray}{HTML}{767676}
\definecolor{linegray}{HTML}{D7DEDE}
\definecolor{softaccent}{HTML}{EAF3F3}

\color{textgray}

% -------------------------------------------------
% Настройки оформления
% -------------------------------------------------
\hypersetup{
  colorlinks=true,
  linkcolor=darkaccent,
  urlcolor=darkaccent
}

\titleformat{\section}
  {\Large\bfseries\color{darkaccent}}
  {\thesection.}{0.55em}{}

\titleformat{\subsection}
  {\large\bfseries\color{accent}}
  {\thesubsection.}{0.5em}{}

\titlespacing*{\section}{0pt}{1.3em}{0.55em}
\titlespacing*{\subsection}{0pt}{1.0em}{0.4em}

\setlist[itemize]{
  leftmargin=1.5em,
  itemsep=0.25em,
  topsep=0.35em
}

\setlist[enumerate]{
  leftmargin=1.8em,
  itemsep=0.55em,
  topsep=0.4em
}

\pagestyle{fancy}
\fancyhf{}
\fancyhead[L]{\small\color{softgray}Путь, перемещение и уравнение координаты}
\fancyhead[R]{\small\color{softgray}Физика • ОГЭ}
\fancyfoot[C]{\small\color{softgray}\thepage}
\renewcommand{\headrulewidth}{0.3pt}
\renewcommand{\headrule}{\hbox to\headwidth{\color{linegray}\leaders\hrule height \headrulewidth\hfill}}

\newcommand{\key}[1]{\textbf{\color{darkaccent}#1}}
\newcommand{\important}[1]{\textbf{#1}}

% -------------------------------------------------
% Документ
% -------------------------------------------------
\begin{document}

% =================================================
% ТИТУЛЬНЫЙ ЛИСТ
% =================================================
\begin{titlepage}
\thispagestyle{empty}

\begin{center}

\vspace*{2.5cm}

{\large\color{softgray}Физика • подготовка к ОГЭ}

\vspace{1.2cm}

{\fontsize{31}{36}\selectfont\bfseries\color{darkaccent}
Чек-лист}

\vspace{0.7cm}

{\LARGE\bfseries
Путь, перемещение\\[0.2em]
и уравнение координаты тела}

\vspace{1.1cm}

{\large
Равномерное прямолинейное движение\\
и работа с координатами}

\vfill

{\large\bfseries
Лёвин Артём Александрович\\[0.3em]
эксклюзивно для Володи Хачатуряна}

\vspace{0.9cm}

{\large \today}

\vspace{1.7cm}

\begin{tikzpicture}[x=1cm,y=1cm]
  \draw[
    color=accent,
    line width=1.1pt
  ]
  (-7,0)
  .. controls (-4.8,0.65) and (-2.8,-0.55) .. (-0.4,0)
  .. controls (2.1,0.55) and (4.6,-0.55) .. (7,0);
\end{tikzpicture}

\vspace{0.7cm}

\end{center}
\end{titlepage}

% =================================================
% ОСНОВНОЙ ЧЕК-ЛИСТ
% =================================================

\section{Главная картина темы}

При движении тела важно различать три связанные, но разные величины:

\begin{itemize}
  \item \key{координата} \(x\) показывает положение тела на выбранной оси;
  \item \key{путь} \(s\) показывает полную длину траектории;
  \item \key{перемещение} показывает изменение положения от начальной точки к конечной.
\end{itemize}

Для движения вдоль оси \(Ox\) изменение координаты записывают так:
\[
\boxed{\Delta x=x-x_0}
\]
где \(x_0\) --- начальная координата, \(x\) --- конечная координата.

Модуль перемещения:
\[
\boxed{|\Delta x|=|x-x_0|}.
\]

\important{Путь и модуль перемещения совпадают только тогда, когда тело двигалось по прямой и не меняло направление.}

\section{Путь и перемещение}

\subsection{Путь}

Путь \(s\) --- длина всей траектории, пройденной телом.

Путь:

\begin{itemize}
  \item является скалярной величиной;
  \item не бывает отрицательным;
  \item зависит от всей траектории движения;
  \item при движении по нескольким участкам равен сумме длин этих участков.
\end{itemize}

Если тело двигалось \(A\to B\to C\), то
\[
s=AB+BC.
\]

\subsection{Перемещение}

Перемещение связывает только начальное и конечное положения тела.

Для движения вдоль оси \(Ox\):
\[
\Delta x=x-x_0.
\]

Знак \(\Delta x\) показывает направление относительно выбранного положительного направления оси.

Модуль перемещения:
\[
|\Delta x|=|x-x_0|.
\]

\subsection{Ключевое различие}

\begin{center}
\renewcommand{\arraystretch}{1.35}
\begin{tabularx}{\linewidth}{@{}p{0.25\linewidth}XX@{}}
\toprule
& \textbf{Путь \(s\)} & \textbf{Перемещение \(\Delta x\)} \\
\midrule
Что учитывает &
всю траекторию &
только начало и конец движения \\

Знак &
\(s\ge 0\) &
может быть положительным, отрицательным или нулевым \\

При возврате назад &
увеличивается &
может уменьшаться по модулю \\

Единица СИ &
метр &
метр \\
\bottomrule
\end{tabularx}
\end{center}

\section{Пример: движение по координатной прямой}

Тело последовательно проходит точки
\[
x_A=-5~\text{м},\qquad
x_B=7~\text{м},\qquad
x_C=2~\text{м}.
\]

\begin{center}
\begin{tikzpicture}[x=0.65cm,y=1cm]
  \draw[->,line width=0.8pt] (-6.3,0) -- (8.2,0) node[right] {\(x\)};
  \foreach \x in {-5,-4,...,7}{
    \draw (\x,0.09)--(\x,-0.09);
  }
  \draw (0,0.13)--(0,-0.13) node[below=5pt] {\(0\)};

  \fill[accent] (-5,0) circle (2pt);
  \fill[accent] (7,0) circle (2pt);
  \fill[accent] (2,0) circle (2pt);

  \node[above=5pt] at (-5,0) {\(A(-5)\)};
  \node[above=5pt] at (7,0) {\(B(7)\)};
  \node[below=7pt] at (2,0) {\(C(2)\)};

  \draw[->,color=accent,line width=1pt,bend left=28]
    (-5,0.25) to node[above] {\(12~\text{м}\)} (7,0.25);

  \draw[->,color=darkaccent,line width=1pt,bend left=30]
    (7,-0.25) to node[below] {\(5~\text{м}\)} (2,-0.25);
\end{tikzpicture}
\end{center}

Путь:
\[
s=AB+BC=|7-(-5)|+|2-7|=12+5=17~\text{м}.
\]

Перемещение:
\[
\Delta x=x_C-x_A=2-(-5)=7~\text{м}.
\]

Модуль перемещения:
\[
|\Delta x|=7~\text{м}.
\]

Итак,
\[
\boxed{s=17~\text{м}},\qquad
\boxed{|\Delta x|=7~\text{м}}.
\]

\section{Движение по дуге окружности}

Если человек прошёл четверть окружности диаметром
\[
d=28~\text{м},
\]
то путь равен длине четверти окружности.

Длина окружности:
\[
L=\pi d.
\]

Следовательно,
\[
s=\frac14\pi d
 =\frac14\pi\cdot 28
 =7\pi~\text{м}.
\]

Радиус:
\[
R=\frac d2=14~\text{м}.
\]

Начальная и конечная точки четверти окружности образуют с центром прямоугольный равнобедренный треугольник.

\begin{center}
\begin{tikzpicture}[scale=0.16]
  \tkzDefPoint(0,0){O}
  \tkzDefPoint(14,0){A}
  \tkzDefPoint(0,14){B}

  \tkzDrawArc[thick,color=accent](O,A)(B)
  \tkzDrawSegments[color=textgray](O,A O,B A,B)

  \tkzDrawPoints(O,A,B)
  \tkzLabelPoints[below](O)
  \tkzLabelPoints[right](A)
  \tkzLabelPoints[above](B)

  \tkzMarkRightAngle[size=2](A,O,B)

  \tkzLabelSegment[below](O,A){\(14\)}
  \tkzLabelSegment[left](O,B){\(14\)}
\end{tikzpicture}
\end{center}

По теореме Пифагора:
\[
AB^2=14^2+14^2=2\cdot14^2.
\]

Поэтому
\[
AB=14\sqrt2~\text{м}.
\]

Получаем:
\[
\boxed{s=7\pi~\text{м}},
\qquad
\boxed{|\Delta\vec r|=14\sqrt2~\text{м}}.
\]

\important{Здесь путь проходит по дуге, а модуль перемещения равен длине хорды между началом и концом движения.}

\section{Скорость}

Для равномерного движения скорость показывает, как быстро изменяется положение тела.

Если движение происходит вдоль оси \(Ox\), удобно использовать проекцию скорости \(v_x\).

Для равномерного прямолинейного движения:
\[
\boxed{v_x=\frac{\Delta x}{\Delta t}}
\]

или
\[
\boxed{v_x=\frac{x-x_0}{t}}
\]
если отсчёт времени начинается при \(t_0=0\).

Обратите внимание:
\[
v_x
\]
может быть отрицательной.

Знак скорости определяет направление движения относительно положительного направления оси.

\begin{itemize}
  \item \(v_x>0\) --- координата увеличивается;
  \item \(v_x<0\) --- координата уменьшается;
  \item \(v_x=0\) --- тело покоится.
\end{itemize}

\important{Отрицательная скорость не означает автоматически движение ``влево''. Она означает движение против положительного направления выбранной оси.}

\section{Путь при равномерном движении}

Если тело движется прямолинейно и не меняет направление, то
\[
s=|v_x|t.
\]

При этом
\[
|\Delta x|=s.
\]

Если направление движения менялось, равенство
\[
s=|\Delta x|
\]
в общем случае уже неверно.

\section{Уравнение координаты}

Основная формула темы:
\[
\boxed{x=x_0+v_xt}.
\]

Она связывает:

\begin{itemize}
  \item \(x\) --- координату тела в момент времени \(t\);
  \item \(x_0\) --- начальную координату;
  \item \(v_x\) --- проекцию скорости на ось \(Ox\);
  \item \(t\) --- время движения.
\end{itemize}

Иногда используют запись
\[
x(t)=x_0+v_xt.
\]

Запись \(x(t)\) означает, что координата рассматривается как функция времени.

\subsection{Как читать формулу}

Формулу
\[
x=x_0+v_xt
\]
полезно проговаривать словами:

\begin{quote}
Конечная координата равна начальной координате плюс изменение координаты, возникающее за время движения.
\end{quote}

\section{Как считывать данные из готового уравнения}

Если дано
\[
x=2+3t,
\]
то
\[
x_0=2,\qquad v_x=3.
\]

Если
\[
x=-1+5t,
\]
то
\[
x_0=-1,\qquad v_x=5.
\]

Если
\[
x=5t,
\]
то отсутствующий свободный член равен нулю:
\[
x_0=0,\qquad v_x=5.
\]

Если
\[
x=5,
\]
то
\[
x_0=5,\qquad v_x=0.
\]

Если
\[
x=-3-10t,
\]
то
\[
x_0=-3,\qquad v_x=-10.
\]

\subsection{Главный алгоритм}

В выражении
\[
x=x_0+v_xt
\]

\begin{itemize}
  \item число \important{без \(t\)} --- начальная координата;
  \item коэффициент \important{при \(t\)} --- проекция скорости.
\end{itemize}

Порядок слагаемых роли не играет:
\[
3t+2=2+3t.
\]

\section{Как найти координату в заданный момент времени}

Пусть
\[
x=15-3t.
\]

Найдём координату при
\[
t=3~\text{с}.
\]

Подставляем значение времени:
\[
x(3)=15-3\cdot3=6.
\]

Если координаты измеряются в метрах, ответ:
\[
\boxed{x=6~\text{м}}.
\]

\important{Единицы координаты должны быть известны из условия задачи. Если условие их не задаёт, самостоятельно приписывать метры нельзя.}

\section{Как составить уравнение координаты}

Пусть
\[
x_0=2~\text{м},
\qquad
v_x=-5~\text{м/с}.
\]

Тогда
\[
x(t)=2-5t.
\]

Через \(10\) секунд:
\[
x(10)=2-5\cdot10=-48~\text{м}.
\]

\section{Как выразить неизвестную величину}

Из
\[
x=x_0+v_xt
\]
можно получить необходимые формулы.

Для скорости:
\[
v_x=\frac{x-x_0}{t},
\qquad t\ne0.
\]

Для времени:
\[
t=\frac{x-x_0}{v_x},
\qquad v_x\ne0.
\]

Для начальной координаты:
\[
x_0=x-v_xt.
\]

\important{Перед делением обязательно проверяйте, что знаменатель не равен нулю.}

\section{Когда лучше сначала подставить числа}

Если преобразование формулы вызывает затруднение, безопасный приём:

\begin{enumerate}
  \item записать исходную физическую формулу;
  \item подставить известные значения;
  \item получить обычное числовое уравнение;
  \item решить его средствами алгебры.
\end{enumerate}

Например:
\[
x=10~\text{м},\qquad
x_0=20~\text{м},\qquad
v_x=-1~\text{м/с}.
\]

Записываем:
\[
10=20-t.
\]

Тогда
\[
t=10~\text{с}.
\]

Такой способ часто снижает риск ошибки при переносе слагаемых.

\section{Физическая проверка ответа}

Результат вычислений необходимо проверить по смыслу.

Пусть
\[
x_0=20~\text{м},\qquad
x=10~\text{м},\qquad
v_x=1~\text{м/с}.
\]

Формально:
\[
10=20+t,
\]
откуда
\[
t=-10~\text{с}.
\]

Для задачи о времени движения после начального момента результат
\[
t<0
\]
не соответствует условию.

Причина видна ещё до вычислений:

\begin{itemize}
  \item начальная координата \(20\);
  \item конечная координата \(10\);
  \item координата должна уменьшиться;
  \item при \(v_x>0\) координата, наоборот, увеличивается.
\end{itemize}

Чтобы задача имела решение при \(t>0\), можно, например, взять
\[
v_x=-1~\text{м/с}.
\]

\section{Движение на координатной плоскости}

Если тело перемещается одновременно по двум координатам, одной оси \(Ox\) уже недостаточно.

Пусть
\[
A(0;2),\qquad B(4;-1).
\]

При прямолинейном движении длина перемещения равна расстоянию между точками:
\[
AB=\sqrt{(4-0)^2+(-1-2)^2}.
\]

Получаем:
\[
AB=\sqrt{16+9}=5.
\]

Если движение из \(A\) в \(B\) заняло \(10~\text{с}\), то модуль скорости при равномерном прямолинейном движении:
\[
v=\frac{AB}{t}=\frac5{10}=0,5.
\]

Единица скорости зависит от единиц координат. Если координаты заданы в метрах, то
\[
v=0,5~\text{м/с}.
\]

\section{Чертёж как инструмент проверки}

Для задач на движение полезно до вычислений сделать схему.

На чертеже желательно показать:

\begin{itemize}
  \item ось и её положительное направление;
  \item начало отсчёта;
  \item начальное положение;
  \item конечное положение;
  \item направление движения;
  \item известные координаты.
\end{itemize}

Чертёж помогает заранее увидеть:

\begin{itemize}
  \item должен ли знак скорости быть положительным или отрицательным;
  \item увеличивается или уменьшается координата;
  \item возможен ли полученный ответ физически;
  \item совпадают ли путь и модуль перемещения.
\end{itemize}

\section{Оформление решения}

Удобный школьный формат:

\[
\text{Дано}\quad\longrightarrow\quad
\text{Найти}\quad\longrightarrow\quad
\text{Решение}\quad\longrightarrow\quad
\text{Ответ}.
\]

Полезный алгоритм:

\begin{enumerate}
  \item Выпишите известные величины с единицами.
  \item Запишите, что требуется найти.
  \item Сделайте схематический чертёж, если он помогает понять движение.
  \item Запишите физическую формулу в общем виде.
  \item Подставьте числа.
  \item Аккуратно решите полученное уравнение.
  \item Проверьте знак и физический смысл.
  \item Укажите единицу полученной величины.
  \item Запишите итоговый ответ.
\end{enumerate}

Форма записи должна помогать проверке и снижать риск вычислительных ошибок. Конкретные требования к оформлению всегда следует сопоставлять с критериями той работы, которую Вы выполняете.

\section{Типичные ошибки}

\begin{enumerate}
  \item \important{Путать путь и перемещение.}

  Если тело возвращалось назад, обычно
  \[
  s>|\Delta x|.
  \]

  \item \important{Писать \(s=x-x_0\) для любого движения.}

  Выражение
  \[
  x-x_0
  \]
  задаёт изменение координаты, а не полный путь.

  \item \important{Терять знак скорости.}

  В уравнении
  \[
  x=-3-10t
  \]
  скорость равна
  \[
  -10,
  \]
  а не \(10\).

  \item \important{Не учитывать скрытый ноль.}

  В
  \[
  x=5t
  \]
  начальная координата равна нулю.

  \item \important{Считать отрицательное время нормальным ответом.}

  В стандартной задаче, где \(t\) --- продолжительность движения после момента старта,
  \[
  t\ge0.
  \]

  \item \important{Делить на ноль при выражении величины.}

  Формула
  \[
  t=\frac{x-x_0}{v_x}
  \]
  применима в таком виде только при
  \[
  v_x\ne0.
  \]

  \item \important{Приписывать единицы, которых нет в условии.}

  Единицы нельзя угадывать.

  \item \important{Решить алгебраическое уравнение с ошибкой после правильной физической модели.}

  Физическая формула и алгебраическое преобразование одинаково важны.
\end{enumerate}

\section{Мини-чек перед завершением задачи}

Перед записью ответа проверьте:

\begin{itemize}
  \item Что дано: путь, перемещение, координата или скорость?
  \item Как направлена ось?
  \item Увеличивается или уменьшается координата?
  \item Верно ли определён знак \(v_x\)?
  \item Не перепутаны ли \(s\) и \(\Delta x\)?
  \item Записана ли исходная формула?
  \item Правильно ли перенесены слагаемые?
  \item Не возникло ли деление на ноль?
  \item Может ли полученное время быть таким по смыслу?
  \item Указана ли корректная единица?
\end{itemize}

% =================================================
% ТРЕНИРОВОЧНЫЙ БЛОК
% =================================================
\clearpage

\section*{Тренировочный блок}
\addcontentsline{toc}{section}{Тренировочный блок}

\begin{center}
{\large\bfseries Путь Б: ежедневное закрепление}\\[0.35em]
{\color{softgray}30.08.2026--04.09.2026 • по 5 заданий в день}
\end{center}

Решайте задачи с полной фиксацией исходной формулы. В задачах на движение вдоль оси желательно выполнять схематический чертёж.

\subsection*{30.08.2026}

\begin{enumerate}[label=\textbf{\arabic*)}]
  \item Тело переместилось из точки \(A\) с координатой
  \[
  x_A=-4~\text{м}
  \]
  в точку \(B\) с координатой \(x_B=9~\text{м}\), а затем в точку \(C\) с координатой \(x_C=3~\text{м}\).
  Найдите путь и модуль перемещения.

  \item Для тела задано уравнение
  \[
  x=6+4t,
  \]
  где \(x\) измеряется в метрах, \(t\) --- в секундах.
  Найдите \(x_0\), \(v_x\) и координату тела через \(5\) секунд.

  \item Начальная координата тела равна
  \[
  x_0=7~\text{м},
  \]
  а скорость
  \[
  v_x=-3~\text{м/с}.
  \]
  Запишите уравнение координаты и найдите \(x\) через \(4\) секунды.

  \item Тело движется равномерно из точки
  \[
  A(1;2)
  \]
  в точку
  \[
  B(7;10)
  \]
  за \(5\) секунд. Координаты заданы в метрах.
  Найдите модуль скорости.

  \item Человек прошёл четверть окружности радиуса \(10~\text{м}\).
  Найдите путь и модуль перемещения.
\end{enumerate}

\subsection*{31.08.2026}

\begin{enumerate}[label=\textbf{\arabic*)}]
  \item По уравнению
  \[
  x=-8+6t
  \]
  определите начальную координату, проекцию скорости и координату при
  \[
  t=3~\text{с}.
  \]

  \item Тело прошло из координаты
  \[
  -7~\text{м}
  \]
  в координату
  \[
  5~\text{м},
  \]
  затем вернулось в координату
  \[
  1~\text{м}.
  \]
  Найдите путь и модуль перемещения.

  \item Известно:
  \[
  x_0=-2~\text{м},\qquad
  v_x=5~\text{м/с},\qquad
  t=7~\text{с}.
  \]
  Найдите конечную координату.

  \item Тело движется по закону
  \[
  x=18-2t.
  \]
  Через сколько секунд тело окажется в точке
  \[
  x=6~\text{м}?
  \]

  \item Тело переместилось из точки
  \[
  A(-2;1)
  \]
  в точку
  \[
  B(4;9)
  \]
  за \(4\) секунды. Координаты заданы в метрах.
  Найдите модуль скорости.
\end{enumerate}

\subsection*{01.09.2026}

\begin{enumerate}[label=\textbf{\arabic*)}]
  \item Уравнение координаты:
  \[
  x=-7t.
  \]
  Найдите \(x_0\), \(v_x\) и координату через \(6\) секунд.

  \item Тело стартовало из точки
  \[
  x_0=12~\text{м}
  \]
  и через \(4\) секунды оказалось в точке
  \[
  x=-8~\text{м}.
  \]
  Найдите \(v_x\).

  \item Тело двигалось по прямой из точки
  \[
  x_A=2~\text{м}
  \]
  в точку
  \[
  x_B=11~\text{м},
  \]
  затем в точку
  \[
  x_C=-1~\text{м}.
  \]
  Найдите путь, перемещение и модуль перемещения.

  \item Точка движется равномерно по закону
  \[
  x=25-5t.
  \]
  Определите момент времени, когда она пересечёт начало координат.

  \item Человек прошёл половину окружности диаметром
  \[
  18~\text{м}.
  \]
  Найдите путь и модуль перемещения.
\end{enumerate}

\subsection*{02.09.2026}

\begin{enumerate}[label=\textbf{\arabic*)}]
  \item Начальная координата:
  \[
  x_0=-10~\text{м},
  \]
  скорость:
  \[
  v_x=4~\text{м/с}.
  \]
  Запишите \(x(t)\) и найдите координату через \(8\) секунд.

  \item Тело за \(6\) секунд переместилось из
  \[
  x_0=14~\text{м}
  \]
  в
  \[
  x=-4~\text{м}.
  \]
  Найдите проекцию скорости.

  \item Уравнение движения:
  \[
  x=9.
  \]
  Определите начальную координату, скорость и координату через \(100\) секунд.

  \item Тело перемещается из
  \[
  A(2;-3)
  \]
  в
  \[
  B(8;5)
  \]
  за \(10\) секунд. Координаты заданы в метрах.
  Найдите длину перемещения и модуль скорости.

  \item Тело проходит точки
  \[
  -3~\text{м}\to 8~\text{м}\to -6~\text{м}.
  \]
  Найдите путь и модуль перемещения.
\end{enumerate}

\subsection*{03.09.2026}

\begin{enumerate}[label=\textbf{\arabic*)}]
  \item Определите \(x_0\) и \(v_x\) из уравнения
  \[
  x=11-4t.
  \]
  Найдите координату при \(t=5~\text{с}\).

  \item Тело движется равномерно с
  \[
  v_x=-6~\text{м/с}
  \]
  и в момент \(t=0\) имеет координату
  \[
  x_0=20~\text{м}.
  \]
  Через сколько секунд координата станет равной
  \[
  -10~\text{м}?
  \]

  \item Начальная координата равна
  \[
  -5~\text{м},
  \]
  конечная координата через \(3\) секунды равна
  \[
  7~\text{м}.
  \]
  Найдите проекцию скорости и запишите уравнение координаты.

  \item Человек прошёл четверть окружности диаметром
  \[
  40~\text{м}.
  \]
  Найдите путь и модуль перемещения.

  \item Тело движется из
  \[
  A(-1;-2)
  \]
  в
  \[
  B(5;6)
  \]
  за \(2\) секунды. Координаты заданы в метрах.
  Найдите модуль скорости.
\end{enumerate}

\subsection*{04.09.2026}

\begin{enumerate}[label=\textbf{\arabic*)}]
  \item Тело проходит последовательно координаты
  \[
  4~\text{м}\to -8~\text{м}\to 1~\text{м}.
  \]
  Найдите путь, перемещение и модуль перемещения.

  \item Уравнение движения:
  \[
  x=-15+3t.
  \]
  Через сколько секунд тело окажется в точке
  \[
  x=12~\text{м}?
  \]

  \item Тело находится в точке
  \[
  x_0=16~\text{м}
  \]
  и движется равномерно. Через \(5\) секунд его координата равна
  \[
  -9~\text{м}.
  \]
  Найдите проекцию скорости и запишите уравнение координаты.

  \item Тело движется по закону
  \[
  x=-4-2t.
  \]
  Определите начальную координату, скорость, направление движения относительно положительного направления \(Ox\) и координату через \(7\) секунд.

  \item Тело перемещается из
  \[
  A(-4;3)
  \]
  в
  \[
  B(8;-2)
  \]
  за \(13\) секунд. Координаты заданы в метрах.
  Найдите длину перемещения и модуль скорости.
\end{enumerate}

% =================================================
% ОТВЕТЫ
% =================================================
\clearpage

\section*{Ответы к тренировочному блоку}
\addcontentsline{toc}{section}{Ответы к тренировочному блоку}

\renewcommand{\arraystretch}{1.35}

\begin{longtable}{@{}p{0.16\linewidth}p{0.08\linewidth}p{0.68\linewidth}@{}}
\toprule
\textbf{Дата} & \textbf{\№} & \textbf{Ответ} \\
\midrule
\endfirsthead

\toprule
\textbf{Дата} & \textbf{\№} & \textbf{Ответ} \\
\midrule
\endhead

\bottomrule
\endfoot

30.08.2026 & 1 &
\(
s=19~\text{м},\quad |\Delta x|=7~\text{м}.
\)
\\

30.08.2026 & 2 &
\(
x_0=6~\text{м},\quad v_x=4~\text{м/с},\quad x(5)=26~\text{м}.
\)
\\

30.08.2026 & 3 &
\(
x=7-3t,\quad x(4)=-5~\text{м}.
\)
\\

30.08.2026 & 4 &
\(
AB=10~\text{м},\quad v=2~\text{м/с}.
\)
\\

30.08.2026 & 5 &
\(
s=5\pi~\text{м},\quad |\Delta\vec r|=10\sqrt2~\text{м}.
\)
\\

\midrule

31.08.2026 & 1 &
\(
x_0=-8~\text{м},\quad v_x=6~\text{м/с},\quad x(3)=10~\text{м}.
\)
\\

31.08.2026 & 2 &
\(
s=16~\text{м},\quad |\Delta x|=8~\text{м}.
\)
\\

31.08.2026 & 3 &
\(
x=33~\text{м}.
\)
\\

31.08.2026 & 4 &
\(
t=6~\text{с}.
\)
\\

31.08.2026 & 5 &
\(
AB=10~\text{м},\quad v=2,5~\text{м/с}.
\)
\\

\midrule

01.09.2026 & 1 &
\(
x_0=0,\quad v_x=-7~\text{м/с},\quad x(6)=-42~\text{м}.
\)
\\

01.09.2026 & 2 &
\(
v_x=-5~\text{м/с}.
\)
\\

01.09.2026 & 3 &
\(
s=21~\text{м},\quad \Delta x=-3~\text{м},\quad |\Delta x|=3~\text{м}.
\)
\\

01.09.2026 & 4 &
\(
t=5~\text{с}.
\)
\\

01.09.2026 & 5 &
\(
s=9\pi~\text{м},\quad |\Delta\vec r|=18~\text{м}.
\)
\\

\midrule

02.09.2026 & 1 &
\(
x=-10+4t,\quad x(8)=22~\text{м}.
\)
\\

02.09.2026 & 2 &
\(
v_x=-3~\text{м/с}.
\)
\\

02.09.2026 & 3 &
\(
x_0=9,\quad v_x=0,\quad x(100)=9.
\)
\\

02.09.2026 & 4 &
\(
AB=10~\text{м},\quad v=1~\text{м/с}.
\)
\\

02.09.2026 & 5 &
\(
s=25~\text{м},\quad |\Delta x|=3~\text{м}.
\)
\\

\midrule

03.09.2026 & 1 &
\(
x_0=11~\text{м},\quad v_x=-4~\text{м/с},\quad x(5)=-9~\text{м}.
\)
\\

03.09.2026 & 2 &
\(
t=5~\text{с}.
\)
\\

03.09.2026 & 3 &
\(
v_x=4~\text{м/с},\quad x=-5+4t.
\)
\\

03.09.2026 & 4 &
\(
s=10\pi~\text{м},\quad |\Delta\vec r|=20\sqrt2~\text{м}.
\)
\\

03.09.2026 & 5 &
\(
AB=10~\text{м},\quad v=5~\text{м/с}.
\)
\\

\midrule

04.09.2026 & 1 &
\(
s=21~\text{м},\quad \Delta x=-3~\text{м},\quad |\Delta x|=3~\text{м}.
\)
\\

04.09.2026 & 2 &
\(
t=9~\text{с}.
\)
\\

04.09.2026 & 3 &
\(
v_x=-5~\text{м/с},\quad x=16-5t.
\)
\\

04.09.2026 & 4 &
\(
x_0=-4~\text{м},\quad
v_x=-2~\text{м/с},\quad
\text{движение против положительного направления }Ox,\quad
x(7)=-18~\text{м}.
\)
\\

04.09.2026 & 5 &
\(
AB=13~\text{м},\quad v=1~\text{м/с}.
\)
\\

\end{longtable}

\end{document}